\frfilename{mt214.tex}
\versiondate{22.5.09}
\copyrightdate{2009}

\def\chaptername{Taksonomi ruang ukur}
\def\sectionname{Subruang}

\newsection{214}

Dalam \S131 saya menjelaskan suatu konstruksi ukuran subruang pada
himpunan bagian terukur.   Kini tiba waktunya memberikan generalisasi
ukuran subruang
pada sembarang himpunan bagian dari suatu ruang ukur.   Hubungan antara
konstruksi ini dan sifat-sifat yang didaftar dalam \S211 tidaklah
selangsung yang mungkin dibayangkan, dan dalam bagian ini saya berusaha
memberikan uraian lengkap tentang apa yang secara umum dapat diharapkan
dari subruang.   Menurut saya, untuk jilid ini hanya (i) subruang umum
dari ruang $\sigma$-hingga dan
(ii) subruang terukur dari ruang ukur umum yang akan benar-benar
diperlukan, dan keduanya tidak menimbulkan kesulitan;  tetapi dalam
jilid-jilid selanjutnya kita akan memerlukan teori lengkapnya.

Saya mulai dengan suatu konstruksi umum untuk `ukuran subruang'
(214A-214C), %214A 214B 214C
beserta uraian tentang integrasi terhadap suatu ukuran subruang
(214E-214G); %214E 214F 214G
hasil-hasil ini (bersama 131E-131H) %131E 131F 131G 131H
memberikan landasan yang kukuh bagi
konsep `integrasi pada suatu himpunan bagian' (214D).   Saya menyajikan
pekerjaan ini dalam keumuman alaminya yang penuh, yang kelak akan mutlak
diperlukan, tetapi bahkan untuk ukuran Lebesgue saja penting untuk
menyadari gagasan-gagasan di sini.    Saya lanjutkan dengan jawaban atas
beberapa pertanyaan yang jelas mengenai ukuran subruang dan sifat-sifat
ruang ukur yang sejauh ini telah dipertimbangkan, baik untuk subruang umum
(214I) maupun subruang terukur (214K), dan saya menyebut suatu konstruksi
dasar untuk menyusun ruang-ruang ukur berdampingan, yakni `jumlah langsung'
dalam 214L-214M. Pada akhir bagian saya membahas suatu masalah perluasan
ukuran (214O-214P).

\leader{214A}{Proposisi} Misalkan $(X,\Sigma,\mu)$ suatu ruang ukur, dan
$Y$ sembarang himpunan bagian dari $X$.   Misalkan $\mu^*$ ukuran luar
yang didefinisikan dari $\mu$\cmmnt{ (132A-132B)}, dan tetapkan
$\Sigma_Y=\{E\cap Y:E\in\Sigma\}$;  misalkan $\mu_Y$ restriksi dari
$\mu^*$ pada $\Sigma_Y$.   Maka $(Y,\Sigma_Y,\mu_Y)$ merupakan ruang ukur.

\proof{{\bf (a)} Telah saya catat dalam 121A bahwa $\Sigma_Y$ merupakan
aljabar-sigma dari himpunan-himpunan bagian $Y$.

\medskip

{\bf (b)} Tentu saja $\mu_YF\in[0,\infty]$ untuk
setiap $F\in\Sigma_Y$.

\medskip

{\bf (c)}  $\mu_Y\emptyset=\mu^*\emptyset=0$.


\medskip

{\bf (d)} Jika $\sequencen{F_n}$ suatu barisan saling lepas dalam
$\Sigma_Y$ dengan gabungan $F$, pilih $E_n$, $E'_n$, $E\in\Sigma$
sedemikian sehingga $F_n=Y\cap E_n$, $F_n\subseteq E'_n$ dan
$\mu_YF_n=\mu E'_n$ untuk setiap $n$, $F\subseteq E$ dan
$\mu_YF=\mu E$ (dengan menggunakan 132Aa berulang kali).   Tetapkan
$G_n=E_n\cap E'_n\cap E\setminus\bigcup_{m<n}E_m$ untuk setiap
$n\in\Bbb N$; maka $\sequencen{G_n}$ saling lepas dan
$F_n\subseteq G_n\subseteq E'_n$ untuk setiap $n$, sehingga
$\mu_YF_n=\mu G_n$.   Selain itu
$F\subseteq\bigcup_{n\in\Bbb N}G_n\subseteq E$, sehingga

\Centerline{$\mu_YF=\mu(\bigcup_{n\in\Bbb N}G_n)
=\sum_{n=0}^{\infty}\mu G_n=\sum_{n=0}^{\infty}\mu_YF_n$.}

\noindent Karena $\sequencen{F_n}$
sembarang, $\mu_Y$ merupakan suatu ukuran.
}%end of proof of 214A

\leader{214B}{Definisi} Jika $(X,\Sigma,\mu)$ sembarang ruang ukur dan
$Y$ sembarang himpunan bagian dari $X$, maka $\mu_Y$, yang didefinisikan
seperti dalam 214A, adalah {\bf ukuran subruang} pada $Y$.

\cmmnt{Patut dicatat hal berikut.}

\leader{214C}{Lema} Misalkan $(X,\Sigma,\mu)$ suatu ruang ukur, $Y$ suatu
himpunan bagian dari $X$, $\mu_Y$ ukuran subruang pada $Y$, dan
$\Sigma_Y$ domainnya.  Maka

(a) untuk setiap $F\in\Sigma_Y$, terdapat $E\in\Sigma$ sedemikian
sehingga $F=E\cap Y$ dan $\mu E=\mu_YF$;

(b) untuk setiap $A\subseteq Y$, $A$ terabaikan terhadap $\mu_Y$ jika dan hanya
jika ia terabaikan terhadap $\mu$;

(c)(i) jika $A\subseteq X$ koterabaikan terhadap $\mu$, maka $A\cap Y$
koterabaikan terhadap $\mu_Y$;

\quad (ii) jika $A\subseteq Y$
koterabaikan terhadap $\mu_Y$, maka $A\cup(X\setminus Y)$ koterabaikan terhadap $\mu$;

(d) $(\mu_Y)^*$\cmmnt{, ukuran luar pada $Y$ yang didefinisikan dari
$\mu_Y$,} berimpit dengan $\mu^*$ pada $\Cal PY$;

(e) jika $Z\subseteq Y\subseteq X$, maka
$\Sigma_Z=(\Sigma_Y)_Z$, yakni aljabar-sigma subruang dari
himpunan-himpunan bagian $Z$ yang dipandang sebagai subruang dari
$(Y,\Sigma_Y)$, dan $\mu_Z=(\mu_Y)_Z$ adalah ukuran subruang pada $Z$
yang dipandang sebagai subruang dari $(Y,\mu_Y)$;

(f) jika $Y\in\Sigma$, maka $\mu_Y$, sebagaimana didefinisikan di sini,
tepat merupakan ukuran subruang pada $Y$ yang didefinisikan dalam
131A-131B\cmmnt{;  yaitu, $\Sigma_Y=\Sigma\cap\Cal PY$ dan
$\mu_Y=\mu\restr\Sigma_Y$}.

\proof{{\bf (a)} Berdasarkan definisi $\Sigma_Y$, terdapat
$E_0\in\Sigma$ sedemikian sehingga $F=E_0\cap Y$.   Berdasarkan 132Aa,
terdapat $E_1\in\Sigma$ sedemikian sehingga $F\subseteq E_1$ dan
$\mu^*F=\mu E_1$.   Tetapkan $E=E_0\cap E_1$;  himpunan ini memenuhi
keperluan.

\medskip

{\bf (b)}(i) Jika $A$ terabaikan terhadap $\mu_Y$, terdapat suatu himpunan
$F\in\Sigma_Y$ sedemikian sehingga $A\subseteq F$ dan $\mu_YF=0$; kini
$\mu^*A\le\mu^*F=0$ sehingga $A$ terabaikan terhadap $\mu$, berdasarkan 132Ad.
(ii) Jika $A$
terabaikan terhadap $\mu$, terdapat $E\in\Sigma$ sedemikian sehingga
$A\subseteq E$ dan $\mu E=0$;  kini $A\subseteq E\cap Y\in\Sigma_Y$ dan
$\mu_Y(E\cap Y)=0$, sehingga $A$ terabaikan terhadap $\mu_Y$.

\medskip

{\bf (c)} Jika $A\subseteq X$ koterabaikan terhadap $\mu$, maka $A\cap Y$
koterabaikan terhadap $\mu_Y$, karena $Y\setminus A=Y\cap(X\setminus A)$
terabaikan terhadap $\mu$, dan karena itu terabaikan terhadap $\mu_Y$.   Jika
$A\subseteq Y$ koterabaikan terhadap $\mu_Y$, maka $A\cup(X\setminus Y)$
koterabaikan terhadap $\mu$ karena $X\setminus(A\cup(X\setminus Y))=Y\setminus A$
terabaikan terhadap $\mu_Y$, dan karena itu terabaikan terhadap $\mu$.

\medskip

{\bf (d)} Misalkan $A\subseteq Y$.   (i) Jika
$A\subseteq E\in\Sigma$, maka $A\subseteq E\cap Y\in\Sigma_Y$, sehingga
$\mu_Y^*A\le\mu_Y(E\cap Y)\le\mu E$;  karena $E$ sembarang,
$\mu_Y^*A\le\mu^*A$.    (ii) Jika $A\subseteq F\in\Sigma_Y$, terdapat
$E\in\Sigma$ sedemikian sehingga $F\subseteq E$ dan
$\mu_YF=\mu^*F=\mu E$; kini $A\subseteq E$ sehingga
$\mu^*A\le\mu E=\mu_YF$.   Karena $F$ sembarang,
$\mu^*A\le\mu_Y^*A$.

\medskip

{\bf (e)} Kesamaan $\Sigma_Z=(\Sigma_Y)_Z$ langsung diperoleh dari
definisi $\Sigma_Y$, dan seterusnya;   kini

\Centerline{$(\mu_Y)_Z=
\mu_Y^*\restr\Sigma_Z=\mu^*\restr\Sigma_Z=\mu_Z$}

\noindent berdasarkan (d).

\medskip

{\bf (f)}  Ini elementer, karena $E\cap Y\in\Sigma$ dan
$\mu^*(E\cap Y)=\mu(E\cap Y)$ untuk setiap $E\in\Sigma$.
}%end of proof of 214C

\leader{214D}{Integrasi pada himpunan bagian:  Definisi} Misalkan
$(X,\Sigma,\mu)$ suatu ruang ukur, $Y$ suatu himpunan bagian dari $X$,
dan $f$ suatu fungsi bernilai $[-\infty,\infty]$ yang didefinisikan pada
suatu himpunan bagian dari $X$.   Dengan $\int_Y f$\cmmnt{ (atau
$\int_Y f(x)\mu(dx)$, dan seterusnya)} saya maksudkan
$\int(f\restrp Y)d\mu_Y$, jika ini ada dalam
$[-\infty,\infty]$\cmmnt{, dengan mengikuti definisi 214A-214B, 133A,
dan 135F, serta mengambil $\dom(f\restrp Y)=Y\cap\dom f$,
$(f\restrp Y)(x)=f(x)$ untuk $x\in Y\cap\dom f$}.   \cmmnt{(Bandingkan
131D.)}

\leader{214E}{Proposisi} Misalkan $(X,\Sigma,\mu)$ suatu ruang ukur,
$Y\subseteq X$, dan $f$ suatu fungsi bernilai $[-\infty,\infty]$ yang
didefinisikan pada suatu himpunan bagian $\dom f$ dari $X$.

(a) Jika $f$ dapat diintegralkan terhadap $\mu$ maka $f\restrp Y$ dapat diintegralkan terhadap $\mu_Y$,
dan $\int_Y f\le\int f$ jika $f$ nonnegatif.

(b) Jika $\dom f\subseteq Y$ dan $f$ dapat diintegralkan terhadap $\mu_Y$, maka terdapat
suatu fungsi yang dapat diintegralkan terhadap $\mu$, yakni $\tilde f$ pada $X$, memperluas
$f$, sedemikian sehingga $\int_F\tilde f=\int_{F\cap Y}f$ untuk setiap
$F\in\Sigma$.

\proof{{\bf (a)(i)} Jika $f$ sederhana terhadap $\mu$, ia dapat dinyatakan sebagai
$\sum_{i=0}^na_i\chi E_i$, dengan $E_0,\ldots,E_n\in\Sigma$,
$a_0,\ldots,a_n\in\Bbb R$ dan $\mu E_i<\infty$ untuk setiap $i$.   Kini
$f\restrp Y=\sum_{i=0}^na_i\chi_Y(E_i\cap Y)$, dengan $\chi_Y(E_i\cap
Y)=(\chi E_i)\restrp Y$ adalah fungsi indikator dari $E_i\cap Y$ yang
dipandang sebagai himpunan bagian dari $Y$;  dan setiap $E_i\cap Y$
termasuk dalam $\Sigma_Y$, dengan $\mu_Y(E_i\cap Y)\le\mu E_i<\infty$,
sehingga $f\restr Y:Y\to\Bbb R$ sederhana terhadap $\mu_Y$.

Jika $f:X\to \Bbb R$ suatu fungsi sederhana nonnegatif, ia dapat
dinyatakan sebagai $\sum_{i=0}^na_i\chi E_i$ dengan
$E_0,\ldots,E_n$ himpunan-himpunan saling lepas berukuran hingga (122Cb).
Kini $f\restrp Y=\sum_{i=0}^na_i\chi_Y(E_i\cap Y)$ dan

\Centerline{$\int (f\restrp Y)d\mu_Y
=\sum_{i=0}^na_i\mu_Y(E_i\cap Y)\le\sum_{i=0}^na_i\mu E_i=\int fd\mu$}

\noindent karena $a_i\ge 0$ setiap kali $E_i\ne\emptyset$, sehingga
$a_i\mu_Y(E_i\cap Y)\le a_i\mu E_i$ untuk setiap $i$.

\medskip

\quad{\bf (ii)} Jika $f$ suatu fungsi nonnegatif yang dapat diintegralkan terhadap $\mu$,
terdapat suatu barisan tak-menurun $\sequencen{f_n}$ dari fungsi-fungsi
sederhana terhadap $\mu$ nonnegatif yang konvergen menuju $f\,\,\mu$-hampir di
mana-mana;  kini $\sequencen{f_n\restrp Y}$ merupakan suatu barisan
tak-menurun dari fungsi-fungsi sederhana terhadap $\mu_Y$ yang meningkat menuju
$f\restrp Y\,\,\mu_Y$-hampir di mana-mana\ (berdasarkan 214Cb), dan

\Centerline{$\sup_{n\in\Bbb N}\int(f_n\restrp Y)d\mu_Y
\le\sup_{n\in\Bbb N}\int f_nd\mu
=\int fd\mu<\infty$,}

\noindent sehingga $\int(f\restrp Y)d\mu_Y$ ada dan tidak lebih dari
$\int fd\mu$.

\medskip

\quad{\bf (iii)} Akhirnya, jika $f$ sembarang fungsi bernilai real yang
dapat diintegralkan terhadap $\mu$, ia dapat dinyatakan sebagai $f_1-f_2$ dengan $f_1$
dan $f_2$ fungsi-fungsi nonnegatif yang dapat diintegralkan terhadap $\mu$, sehingga
$f\restrp Y=(f_1\restrp Y)-(f_2\restrp Y)$ dapat diintegralkan terhadap $\mu_Y$.

\medskip

{\bf (b)} Kita katakan bahwa jika $f$ suatu fungsi dapat diintegralkan terhadap $\mu_Y$,
maka suatu `perluasan penyelubung' dari $f$ adalah fungsi
dapat diintegralkan terhadap $\mu$, $\tilde f$, yang memperluas $f$, bernilai real pada
$X\setminus Y$, sedemikian sehingga $\int_F\tilde f=\int_{F\cap Y}f$
untuk setiap $F\in\Sigma$.

\medskip

\quad{\bf (i)} Jika $f$ berbentuk $\chi H$, dengan $H\in\Sigma_Y$ dan
$\mu_YH<\infty$, ambil $E_0\in\Sigma$ sedemikian sehingga $H=Y\cap E_0$
dan $E_1\in\Sigma$ suatu selubung terukur bagi $H$ (132Ee);   maka
$E=E_0\cap E_1$ merupakan selubung terukur bagi $H$ dan $H=E\cap Y$.
Tetapkan $\tilde f=\chi E$, yang dipandang sebagai fungsi dari $X$ ke
$\{0,1\}$.   Maka $\tilde f\restrp Y=f$, dan untuk setiap $F\in\Sigma$
kita mempunyai

\Centerline{$\int_F\tilde f=\mu_F(E\cap F)=\mu(E\cap F)
=\mu^*(H\cap F)=\mu_{Y\cap F}(H\cap F)=\int_{Y\cap F}f$.}

\noindent Jadi $\tilde f$ merupakan perluasan penyelubung dari $f$.

\medskip

\quad{\bf (ii)} Jika $f$, $g$ fungsi-fungsi dapat diintegralkan terhadap $\mu_Y$
dengan perluasan penyelubung $\tilde f$, $\tilde g$, dan $a$,
$b\in\Bbb R$, maka $a\tilde f+b\tilde g$ memperluas $af+bg$ dan

$$\eqalign{\int_Fa\tilde f+b\tilde g
&=a\int_F\tilde f+b\int_F\tilde g\cr
&=a\int_{F\cap Y}f+b\int_{F\cap Y}g
=\int_{F\cap Y}af+bg\cr}$$

\noindent untuk setiap $F\in\Sigma$, sehingga $a\tilde f+b\tilde g$
merupakan perluasan penyelubung dari $af+bg$.

\medskip

\quad{\bf (iii)} Dengan menggabungkan (i) dan (ii), kita melihat bahwa
setiap fungsi sederhana terhadap $\mu_Y$, $f$, mempunyai suatu perluasan
penyelubung.

\medskip

\quad{\bf (iv)} Sekarang andaikan $\sequencen{f_n}$ suatu barisan
tak-menurun dari fungsi-fungsi sederhana terhadap $\mu_Y$ nonnegatif yang konvergen
$\mu_Y$-hampir di mana-mana menuju suatu fungsi yang
dapat diintegralkan terhadap $\mu_Y$, yakni $f$.   Untuk setiap $n\in\Bbb N$ misalkan
$\tilde f_n$ suatu perluasan penyelubung dari $f_n$. Maka
$\tilde f_n\leae\tilde f_{n+1}$.   \Prf\ Jika $F\in\Sigma$ maka

\Centerline{$\int_F\tilde f_n=\int_{F\cap Y}f_n
\le\int_{F\cap Y}f_{n+1}=\int_F\tilde f_{n+1}$.}

\noindent Jadi $\tilde f_n\leae\tilde f_{n+1}$, berdasarkan 131Ha.\
\QeD\   Selain itu

\Centerline{$\lim_{n\to\infty}\int_F\tilde f_n
=\lim_{n\to\infty}\int_{F\cap Y}f_n
=\int_{F\cap Y}f$}

\noindent untuk setiap $F\in\Sigma$.   Dengan terlebih dahulu mengambil
$F=X$, teorema B. Levi memberi tahu kita bahwa
$h=\lim_{n\to\infty}\tilde f_n$ terdefinisi (sebagai fungsi bernilai
real) $\mu$-hampir di mana-mana;  kini dengan membiarkan $F$ bervariasi,
kita mempunyai $\int_Fh=\int_{F\cap Y}f$ untuk setiap $F\in\Sigma$,
karena $h\restr F=\lim_{n\to\infty}\tilde f_n\restr F\,\,\mu_F$-hampir di mana-mana.
(Tampaknya saya menggunakan 214Cb di sini.)   Kini
$h\restrp Y=f\,\,\mu_Y$-hampir di mana-mana, sekali lagi berdasarkan 214Cb.   Jika kita
mendefinisikan $\tilde f$ dengan menetapkan

\Centerline{$\tilde f(x)=f(x)$ untuk $x\in\dom f$, $h(x)$ untuk
$x\in\dom h\setminus\dom f$, $0$ untuk $x\in X$ lainnya,}

\noindent maka $\tilde f$ terdefinisi di seluruh $X$ dan sama dengan
$h\,\,\mu$-hampir di mana-mana;  sehingga jika $F\in\Sigma$,
$\tilde f\restr F$ akan sama dengan $h\restr F\,\,\mu_F$-hampir di
mana-mana, dan

\Centerline{$\int_F\tilde f=\int_Fh=\int_{F\cap Y}f$.}

\noindent Karena $F$ sembarang, $\tilde f$ merupakan perluasan
penyelubung dari $f$.

\medskip

\quad{\bf (v)} Jadi setiap fungsi dapat diintegralkan terhadap $\mu_Y$ nonnegatif
mempunyai suatu perluasan penyelubung.   Dengan kembali menggunakan (ii),
setiap fungsi dapat diintegralkan terhadap $\mu_Y$ mempunyai suatu perluasan
penyelubung, sebagaimana dinyatakan.
}%end of proof of 214E

\leader{214F}{Proposisi} Misalkan $(X,\Sigma,\mu)$ suatu ruang ukur, $Y$
suatu himpunan bagian dari $X$, dan $f$ suatu fungsi bernilai
$[-\infty,\infty]$ sedemikian sehingga $\int_Xf$ terdefinisi dalam
$[-\infty,\infty]$.   Jika {\it salah satu} dari hal berikut berlaku: $Y$
mempunyai ukuran luar penuh dalam $X$ {\it atau} $f$ bernilai nol hampir
di mana-mana pada $X\setminus Y$, maka $\int_Yf$ terdefinisi dan sama
dengan $\int_Xf$.

\proof{{\bf (a)} Mula-mula andaikan $f$ nonnegatif,
terukur terhadap $\Sigma$, dan terdefinisi di seluruh $X$.   Dalam hal ini
$f\restrp Y$ terukur terhadap $\Sigma_Y$.   Tetapkan
$F_{nk}=\{x:x\in X,\,f(x)\ge 2^{-n}k\}$ untuk $k$, $n\in\Bbb N$,
$f_n=\sum_{k=1}^{4^n}2^{-n}\chi F_{nk}$ untuk $n\in\Bbb N$, sehingga
$\sequencen{f_n}$ merupakan barisan tak-menurun dari fungsi-fungsi
terukur bernilai real yang konvergen di setiap titik menuju $f$, dan
$\int_Xf=\lim_{n\to\infty}\int_Xf_n$.   Untuk setiap $n\in\Bbb N$ dan $k\ge 1$,

\Centerline{$\mu_Y(F_{nk}\cap Y)=\mu^*(F_{nk}\cap Y)=\mu F_{nk}$}

\noindent entah karena $F_{nk}\setminus Y$ terabaikan atau karena
$X$ merupakan selubung terukur bagi $Y$.   Jadi

$$\eqalignno{\int_Yf
&=\lim_{n\to\infty}\int_Yf_n
=\lim_{n\to\infty}\sum_{k=1}^{4^n}2^{-n}\mu_Y(F_{nk}\cap Y)\cr
&=\lim_{n\to\infty}\sum_{k=1}^{4^n}2^{-n}\mu F_{nk}
=\lim_{n\to\infty}\int_Xf_n
=\int_Xf.\cr}$$

\medskip

{\bf (b)} Sekarang andaikan $f$ nonnegatif, terdefinisi hampir
di mana-mana dalam $X$, dan terukur secara virtual terhadap $\mu$.
Dalam hal ini terdapat suatu himpunan terukur $E\subseteq\dom f$ yang
koterabaikan terhadap $\mu$ sedemikian sehingga $f\restr E$ terukur. Tetapkan
$\tilde f(x)=f(x)$ untuk $x\in E$, $0$ untuk $x\in X\setminus E$; maka
$\tilde f$ memenuhi syarat-syarat (a) dan $f=\tilde f\,\,\mu$-hampir di mana-mana.
Oleh karena itu $f\restrp Y=\tilde f\restrp Y\,\,\mu_Y$-hampir di mana-mana\ (214Cc), dan

\Centerline{$\int_Yf=\int_Y\tilde f=\int_X\tilde f=\int_Xf$.}

\medskip

{\bf (c)} Akhirnya, untuk kasus umum, kita dapat menerapkan (b) pada
bagian positif dan negatif $f^+$, $f^-$ dari $f$ untuk memperoleh

\Centerline{$\int_Yf=\int_Yf^+-\int_Yf^-=\int_Xf^+-\int_Xf^-=\int_Xf$.}
}%end of proof of 214F

\leader{214G}{Korolari} Misalkan $(X,\Sigma,\mu)$ suatu ruang ukur, $Y$
suatu himpunan bagian dari $X$, dan $E\in\Sigma$ suatu selubung terukur
bagi $Y$.   Jika $f$ suatu fungsi bernilai $[-\infty,\infty]$ sedemikian
sehingga $\int_Ef$ terdefinisi dalam $[-\infty,\infty]$, maka $\int_Yf$
terdefinisi dan sama dengan $\int_Ef$.

\proof{ Berdasarkan 214Ce, kita dapat mengidentifikasi ukuran subruang
$\mu_Y$ dengan ukuran subruang $(\mu_E)_Y$ yang diinduksi oleh ukuran
subruang pada $E$.   Kini, jika dipandang sebagai subruang dari $E$, $Y$
mempunyai ukuran luar penuh, sehingga 214F memberikan hasilnya.
}%end of proof of 214G

\leader{214H}{Subruang dan \dvrocolon{metode} \Caratheodory}\cmmnt{
Hasil-hasil teknis mudah berikut kadang-kadang akan berguna.

\medskip

\noindent}{\bf Lema} Misalkan $X$ suatu himpunan, $Y\subseteq X$ suatu
himpunan bagian, dan $\theta$ suatu ukuran luar pada $X$.

(a) $\theta_Y=\theta\restrp\Cal PY$ merupakan suatu ukuran luar pada $Y$.

(b) Misalkan $\mu$, $\nu$ berturut-turut merupakan ukuran pada $X$, $Y$
yang didefinisikan dengan metode \Caratheodory\ dari ukuran-ukuran luar
$\theta$, $\theta_Y$, dan $\Sigma$, $\Tau$ domainnya; misalkan $\mu_Y$
ukuran subruang pada $Y$ yang diinduksi oleh $\mu$, dan $\Sigma_Y$
domainnya.   Maka

\quad (i) $\Sigma_Y\subseteq\Tau$ dan $\nu F\le\mu_YF$ untuk setiap
$F\in\Sigma_Y$;

\quad(ii) jika $Y\in\Sigma$ maka $\nu=\mu_Y$;

\quad(iii) jika $\theta=\mu^*$\cmmnt{ (yakni, $\theta$ bersifat
`reguler')} maka $\nu$ memperluas $\mu_Y$;

\quad(iv) jika $\theta=\mu^*$ dan $\theta Y<\infty$ maka $\nu=\mu_Y$.

\proof{{\bf (a)} Anda hanya perlu membaca definisi `ukuran luar' (113A).

\medskip

{\bf (b)(i)} Andaikan $F\in\Sigma_Y$.   Maka himpunan ini berbentuk $E\cap Y$
dengan $E\in\Sigma$.   Jika $A\subseteq Y$, maka

\Centerline{$\theta_Y(A\cap F)+\theta_Y(A\setminus F)
=\theta(A\cap F)+\theta(A\setminus F)
=\theta(A\cap E)+\theta(A\setminus E)
=\theta A=\theta_YA$,}

\noindent sehingga $F\in\Tau$.   Kini

\Centerline{$\nu F=\theta_YF=\theta F\le\mu^*F=\mu_YF$.}

\medskip

\quad{\bf (ii)} Andaikan $F\in\Tau$.   Jika $A\subseteq X$, maka

$$\eqalign{\theta A
&=\theta(A\cap Y)+\theta(A\setminus Y)
=\theta_Y(A\cap Y)+\theta(A\setminus Y)\cr
&=\theta_Y(A\cap Y\cap F)+\theta_Y(A\cap Y\setminus F)
  +\theta(A\setminus Y)\cr
&=\theta(A\cap F)+\theta(A\cap Y\setminus F)+\theta(A\setminus Y)\cr
&=\theta(A\cap F)+\theta((A\setminus F)\cap Y)
  +\theta((A\setminus F)\setminus Y)
=\theta(A\cap F)+\theta(A\setminus F);\cr}$$

\noindent karena $A$ sembarang, $F\in\Sigma$ dan karena itu
$F\in\Sigma_Y$.   Selain itu

\Centerline{$\mu_YF=\mu F=\theta F=\theta_YF=\nu F$.}
\noindent Dengan menggabungkan hal ini dengan (i), kita melihat bahwa
$\mu_Y$ dan $\nu$ identik.

\medskip

\quad{\bf (iii)} Misalkan $F\in\Sigma_Y$.   Maka $F\in\Tau$, berdasarkan
(i). Kini $\nu F=\theta F=\mu^*F=\mu_YF$.   Karena $F$ sembarang,
$\nu$ memperluas $\mu_Y$.

\medskip

\quad{\bf (iv)} Sekarang andaikan $F\in\Tau$.   Karena
$\mu^*Y=\theta Y<\infty$, terdapat selubung-selubung terukur $E_1$, $E_2$
bagi $F$ dan $Y\setminus F$ terhadap $\mu$ (132Ee).   Maka

$$\eqalign{\theta Y
&=\theta_YY
=\theta_YF+\theta_Y(Y\setminus F)
=\theta F+\theta(Y\setminus F)\cr
&=\mu^*F+\mu^*(Y\setminus F)
=\mu E_1+\mu E_2\ge\mu(E_1\cup E_2)
=\theta(E_1\cup E_2)\ge\theta Y,\cr}$$

\noindent sehingga $\mu E_1+\mu E_2=\mu(E_1\cup E_2)$ dan

\Centerline{$\mu(E_1\cap E_2)=\mu E_1+\mu E_2-\mu(E_1\cup E_2)=0$.}

\noindent Karena $\mu$ lengkap (212A) dan
$E_1\cap Y\setminus F\subseteq E_1\cap E_2$ terabaikan terhadap $\mu$,
himpunan tersebut termasuk dalam $\Sigma$; maka
$F=Y\cap(E_1\setminus(E_1\cap Y\setminus F))$ termasuk dalam
$\Sigma_Y$.   Jadi $\Tau\subseteq\Sigma_Y$; dengan menggabungkan hal ini
dengan (iii), kita melihat bahwa $\nu=\mu_Y$.
}%end of proof of 214H

\leader{214I}{}\cmmnt{ Sekarang saya beralih ke hubungan antara
ukuran subruang dan klasifikasi ruang ukur yang dikembangkan dalam bab ini.

\medskip

\noindent}{\bf Teorema} Misalkan $(X,\Sigma,\mu)$ suatu ruang ukur dan
$Y$ suatu himpunan bagian dari $X$.   Misalkan $\mu_Y$ ukuran subruang
pada $Y$ dan $\Sigma_Y$ domainnya.

(a) Jika $(X,\Sigma,\mu)$ lengkap, atau berhingga total, atau
$\sigma$-hingga, atau dapat dilokalkan secara ketat, maka
$(Y,\Sigma_Y,\mu_Y)$ juga demikian. Jika
$\langle X_i\rangle_{i\in I}$ suatu dekomposisi $X$ untuk $\mu$, maka
$\langle X_i\cap Y\rangle_{i\in I}$ suatu dekomposisi $Y$ untuk $\mu_Y$.

(b) Dengan menuliskan $\hat\mu$ untuk pelengkapan $\mu$, ukuran subruang
$\hat\mu_Y=(\hat\mu)_Y$ merupakan pelengkapan $\mu_Y$.

(c) Jika $(X,\Sigma,\mu)$ mempunyai himpunan-himpunan terabaikan yang
ditentukan secara lokal, maka $\mu_Y$ semihingga.

(d) Jika $(X,\Sigma,\mu)$ lengkap dan ditentukan secara lokal, maka
$(Y,\Sigma_Y,\mu_Y)$ lengkap dan semihingga.

(e) Jika $(X,\Sigma,\mu)$ lengkap, ditentukan secara lokal, dan dapat
dapat dilokalkan, maka $(Y,\Sigma_Y,\mu_Y)$ juga demikian.

\proof{{\bf (a)(i)} Andaikan $(X,\Sigma,\mu)$ lengkap. Jika
$A\subseteq U\in\Sigma_Y$ dan $\mu_YU=0$, terdapat $E\in\Sigma$
sedemikian sehingga $U=E\cap Y$ dan $\mu E=\mu_YU=0$; kini
$A\subseteq E$, sehingga $A\in\Sigma$ dan $A=A\cap Y\in\Sigma_Y$.

\medskip

\quad{\bf (ii)}
$\mu_YY=\mu^*Y\le\mu X$, sehingga $\mu_Y$ berhingga total jika $\mu$
berhingga total.

\medskip

\quad{\bf (iii)} Jika $\sequencen{X_n}$ suatu barisan himpunan
berukuran hingga terhadap $\mu$ yang menutupi $X$, maka
$\sequencen{X_n\cap Y}$ suatu barisan himpunan berukuran hingga terhadap
$\mu_Y$ yang menutupi $Y$.   Jadi $(Y,\Sigma_Y,\mu_Y)$ bersifat
$\sigma$-hingga jika $(X,\Sigma,\mu)$ bersifat demikian.

\medskip

\quad{\bf (iv)} Andaikan $\langle X_i\rangle_{i\in I}$ suatu
dekomposisi $X$ untuk $\mu$.   Maka $\langle X_i\cap
Y\rangle_{i\in I}$ merupakan dekomposisi $Y$ untuk $\mu_Y$.   \Prf\
Karena $\mu_Y(X_i\cap Y)\le\mu X_i<\infty$ untuk setiap $i$,
$\langle X_i\cap Y\rangle_{i\in I}$ merupakan partisi $Y$ menjadi
himpunan-himpunan berukuran hingga.   Andaikan $U\subseteq Y$ sedemikian
sehingga $U_i=U\cap X_i\cap Y\in\Sigma_Y$ untuk setiap $i$.   Untuk setiap
$i\in I$, pilih $E_i\in\Sigma$ sedemikian sehingga $U_i=E_i\cap Y$ dan
$\mu E_i=\mu_YU_i$; tentu saja kita dapat mengandaikan
$E_i\subseteq X_i$.   Tetapkan $E=\bigcup_{i\in I}E_i$.   Maka
$E\cap X_i=E_i\in\Sigma$ untuk setiap $i$, sehingga $E\in\Sigma$ dan
$\mu E=\sum_{i\in I}\mu E_i$.   Kini $U=E\cap Y$, sehingga
$U\in\Sigma_Y$ dan

\Centerline{$\mu_YU\le\mu E=\sum_{i\in I}\mu E_i
=\sum_{i\in I}\mu_YU_i$.}

\noindent Di pihak lain, $\mu_YU$ tentu lebih besar dari atau sama dengan
$\sum_{i\in I}\mu_YU_i
=\sup_{J\subseteq I\text{ berhingga}}\sum_{i\in J}\mu_YU_i$, sehingga
keduanya sama.   Karena $U$ sembarang,
$\langle X_i\cap Y\rangle_{i\in I}$ merupakan dekomposisi $Y$ untuk
$\mu_Y$.\ \Qed

Akibatnya $(Y,\Sigma_Y,\mu_Y)$ dapat dilokalkan secara ketat jika
$(X,\Sigma,\mu)$ dapat dilokalkan secara ketat.

\medskip

{\bf (b)} Domain pelengkapan $(\mu_Y)\sphat\mskip5mu$ adalah

$$\eqalignno{\hat\Sigma_Y
&=\{F\symmdiff A:F\in\Sigma_Y,\,A\subseteq Y
   \text{ bersifat }\mu_Y\text{-terabaikan}\}\cr
&=\{(E\cap Y)\symmdiff(A\cap Y):E\in\Sigma,\,
   A\subseteq X\text{ bersifat }\mu\text{-terabaikan}\}\cr
\displaycause{214Cb}
&=\{(E\symmdiff A)\cap Y:E\in\Sigma,\,
   A\text{ bersifat }\mu\text{-terabaikan}\}
=\dom\hat\mu_Y.\cr}$$

\noindent Jika $H\in\hat\Sigma_Y$, maka

\Centerline{$(\mu_Y)\sphat\mskip3mu(H)
=\mu_Y^*H=\mu^*H=(\hat\mu)^*H=\hat\mu_YH$,}

\noindent dengan menggunakan 214Cd untuk langkah kedua dan 212Ea untuk
langkah ketiga.

\medskip

{\bf (c)} Ambil $U\in\Sigma_Y$ sedemikian sehingga $\mu_YU>0$.   Maka
terdapat $E\in\Sigma$ sedemikian sehingga $\mu E<\infty$ dan
$\mu^*(E\cap U)>0$.   \Prf\Quer\ Jika tidak, $E\cap U$ terabaikan
terhadap $\mu$ setiap kali $\mu E<\infty$; karena $\mu$ mempunyai
himpunan terabaikan yang ditentukan secara lokal, $U$ terabaikan terhadap
$\mu$ dan $\mu_YU=\mu^*U=0$.\ \Bang\QeD   Kini
$E\cap U\in\Sigma_Y$ dan

\Centerline{$0<\mu^*(E\cap U)=\mu_Y(E\cap U)\le\mu E<\infty$.}

\medskip

{\bf (d)} Berdasarkan (a), $\mu_Y$ lengkap; berdasarkan 213J dan (c) di
sini, ukuran tersebut semihingga.

\medskip

{\bf (e)} Berdasarkan (d), $\mu_Y$ lengkap dan semihingga. Untuk melihat
bahwa ukuran tersebut ditentukan secara lokal, ambil sembarang $U\subseteq Y$
sedemikian sehingga $U\cap V\in\Sigma_Y$ setiap kali $V\in\Sigma_Y$ dan
$\mu_YV<\infty$.   Berdasarkan 213J dan 213L, terdapat selubung terukur
$E$ bagi $U$ terhadap $\mu$; tentu saja $E\cap Y\in\Sigma_Y$.

Saya menyatakan bahwa $\mu(E\cap Y\setminus U)=0$.   \Prf\ Ambil
sembarang $F\in\Sigma$ dengan $\mu F<\infty$.   Maka
$F\cap U\in\Sigma_Y$, sehingga
\Centerline{$\mu_Y(F\cap E\cap Y)\le\mu(F\cap E)=\mu^*(F\cap
U)=\mu_Y(F\cap U)\le\mu_Y(F\cap E\cap Y)$;}

\noindent dengan demikian $\mu_Y(F\cap E\cap Y)=\mu_Y(F\cap U)$ dan

\Centerline{$\mu^*(F\cap E\cap Y\setminus U)
=\mu_Y(F\cap E\cap Y\setminus U)=0$.}

\noindent Karena $\mu$ lengkap, $\mu(F\cap E\cap Y\setminus U)=0$;
karena $\mu$ ditentukan secara lokal dan $F$ sembarang,
$\mu(E\cap Y\setminus U)=0$.\ \QeD\  Tetapi ini berarti
$E\cap Y\setminus U\in\Sigma_Y$ dan $U\in\Sigma_Y$.   Karena $U$
sembarang, $\mu_Y$ ditentukan secara lokal.

Untuk melihat bahwa $\mu_Y$ dapat dilokalkan, misalkan $\Cal U$
sembarang keluarga dalam $\Sigma_Y$.   Tetapkan

\Centerline{$\Cal E=\{E:E\in\Sigma,\,\mu E<\infty,\,
  \mu E=\mu^*(E\cap U)$ untuk suatu $U\in\Cal U\}$,}

\noindent dan misalkan $G\in\Sigma$ suatu supremum esensial bagi
$\Cal E$ dalam $\Sigma$.   Saya menyatakan bahwa $G\cap Y$ merupakan
supremum esensial bagi $\Cal U$ dalam $\Sigma_Y$.   \Prf\ (i) \Quer\
Jika $U\in\Cal U$ dan $U\setminus(G\cap Y)$ tidak terabaikan, maka
(karena $\mu_Y$ semihingga) terdapat $V\in\Sigma_Y$ sedemikian sehingga
$V\subseteq U\setminus G$ dan $0<\mu_YV<\infty$.   Kini terdapat
$E\in\Sigma$ sedemikian sehingga $V\subseteq E$ dan $\mu E=\mu^*V$.
Kita mempunyai $\mu^*(E\cap U)\ge\mu^*V=\mu E$, sehingga $E\in\Cal E$
dan $E\setminus G$ harus terabaikan; tetapi
$V\subseteq E\setminus G$ tidak terabaikan.\ \BanG\   Jadi
$U\setminus(G\cap Y)$ terabaikan untuk setiap $U\in\Cal U$.
(ii) Jika $W\in\Sigma_Y$ sedemikian sehingga $U\setminus W$ terabaikan
untuk setiap $U\in\Cal U$, nyatakan $W$ sebagai $H\cap Y$ dengan
$H\in\Sigma$.   Jika $E\in\Cal E$, terdapat $U\in\Cal U$ sedemikian
sehingga $\mu E=\mu^*(E\cap U)$; kini $\mu^*(E\cap U\setminus W)=0$,
sehingga $\mu E=\mu^*(E\cap U\cap W)\le\mu(E\cap H)$ dan
$E\setminus H$ terabaikan.   Karena $E$ sembarang, $H$ merupakan batas
atas esensial bagi $\Cal E$ dan $G\setminus H$ terabaikan; tetapi ini
berarti $G\cap Y\setminus W$ terabaikan.   Karena $W$ sembarang,
$G\cap Y$ merupakan supremum esensial bagi $\Cal U$.\ \Qed

Karena $\Cal U$ sembarang, $\mu_Y$ dapat dilokalkan.
}%end of proof of 214I

\leader{214J}{Integral \dvrocolon{atas dan bawah}}\cmmnt{ Fakta-fakta
elementer berikut kadang-kadang berguna.

\medskip

\noindent}{\bf Proposisi} Misalkan $(X,\Sigma,\mu)$ suatu ruang ukur,
$A$ suatu himpunan bagian dari $X$, dan $f$ suatu fungsi bernilai real
yang terdefinisi hampir di mana-mana dalam $X$.   Maka

(a) jika {\it salah satu} dari hal berikut berlaku: $f$ nonnegatif
{\it atau} $A$ mempunyai ukuran luar penuh dalam $X$,
$\overline{\intop}(f\restr A)d\mu_A\le\overline{\intop}fd\mu$;

(b) jika $A$ mempunyai ukuran luar penuh dalam $X$,
$\underline{\int}fd\mu\le\underline{\int}(f\restr A)d\mu_A$.

\proof{{\bf (a)(i)} Andaikan $f$ nonnegatif. Jika
$\overline{\intop}fd\mu=\infty$, hasilnya sepele. Jika tidak, berdasarkan
133J(a-i) terdapat suatu fungsi yang terintegralkan terhadap $\mu$, yakni $g$,
sedemikian sehingga $f\le g\,\,\mu$-hampir di mana-mana\ dan
$\overline{\intop}fd\mu=\int g\,d\mu$.   Kini
$f\restr A\le g\restr A\,\,\mu_A$-hampir di mana-mana, berdasarkan 214Cb, dan
$\int(g\restr A)\,d\mu_A$ terdefinisi serta tidak lebih dari
$\int g\,d\mu$, berdasarkan 214Ea; jadi

\Centerline{$\overlineint(f\restr A)d\mu_A\le\int(g\restr A)d\mu_A
\le\int g\,d\mu=\overlineint fd\mu$.}

\medskip

\quad{\bf (ii)} Sekarang andaikan $A$ mempunyai ukuran luar penuh dalam
$X$. Jika $g$ sedemikian sehingga $f\le g\,\,\mu$-hampir di mana-mana\ dan
$\int g\,d\mu$ terdefinisi dalam $[-\infty,\infty]$, maka
$f\restr A\le g\restr A\,\,\mu_A$-hampir di mana-mana\ dan
$\int(g\restr A)d\mu_A=\int g\,d\mu$, berdasarkan 214F.   Jadi
$\overline{\int}(f\restr A)d\mu_A\le\int g\,d\mu$.   Karena $g$
sembarang, $\overline{\int}(f\restr A)d\mu_A\le\overline{\int}f\,d\mu$.

\medskip

{\bf (b)} Terapkan (a) pada $-f$, dan gunakan 133J(b-iv).
}%end of proof of 214J

\leader{214K}{Subruang terukur:  Proposisi}
Misalkan $(X,\Sigma,\mu)$ suatu ruang ukur.

(a) Misalkan $E\in\Sigma$ dan $\mu_E$ ukuran subruang, dengan
$\Sigma_E$ sebagai domainnya.   Jika $(X,\Sigma,\mu)$ lengkap, atau
berhingga total, atau $\sigma$-hingga, atau dapat dilokalkan secara
ketat, atau semihingga, atau dapat dilokalkan, atau ditentukan secara
lokal, atau tanpa atom, atau atomik murni, maka
$(E,\Sigma_E,\mu_E)$ juga demikian.

(b) Andaikan $\langle X_i\rangle_{i\in I}$ suatu partisi $X$ menjadi
himpunan-himpunan terukur\cmmnt{ (yang tidak harus berukuran hingga)}
sedemikian sehingga

\Centerline{$\Sigma=\{E:E\subseteq X,\,
E\cap  X_i\in\Sigma$ untuk setiap $i\in I\}$,}

\Centerline{$\mu E=\sum_{i\in I}\mu(E\cap X_i)$ untuk setiap $E\in\Sigma$.}
\noindent Maka $(X,\Sigma,\mu)$ lengkap, atau dapat dilokalkan secara
ketat, atau semihingga, atau dapat dilokalkan, atau ditentukan secara
lokal, atau tanpa atom, atau atomik murni, jika dan hanya jika
$(X_i,\Sigma_{X_i},\mu_{X_i})$ mempunyai sifat tersebut untuk setiap
$i\in I$.

\proof{ Saya sungguh berpendapat bahwa jika Anda telah membaca dengan
saksama sampai titik ini, Anda semestinya menganggap hal ini mudah.
Jika masih ragu, ini merupakan kumpulan enam belas latihan yang sangat
cocok untuk dikerjakan.
}%end of proof of 214K

\leader{214L}{Jumlah langsung} Misalkan
$\langle(X_i,\Sigma_i,\mu_i)\rangle_{i\in I}$ sembarang keluarga
berindeks dari ruang-ruang ukur.   Tetapkan
$X=\bigcup_{i\in I}(X_i\times\{i\})$; untuk $E\subseteq X$, $i\in I$
tetapkan $E_i=\{x:(x,i)\in E\}$.   Tuliskan

\Centerline{$\Sigma
=\{E:E\subseteq X,\,E_i\in\Sigma_i$ untuk setiap $i \in I\}$,}

\Centerline{$\mu E=\sum_{i\in I}\mu_iE_i$ untuk setiap $E\in\Sigma$.}

\noindent Maka mudah diperiksa bahwa $(X,\Sigma,\mu)$ merupakan ruang
ukur; saya akan menyebutnya {\bf jumlah langsung} dari keluarga
$\familyiI{(X_i,\Sigma_i,\mu_i)}$.   Perhatikan bahwa jika
$(X,\Sigma,\mu)$ sembarang ruang ukur yang dapat dilokalkan secara
ketat, dengan dekomposisi $\familyiI{X_i}$, maka kita mempunyai isomorfisme
alami antara $(X,\Sigma,\mu)$ dan jumlah langsung
$(X',\Sigma',\mu')=\bigoplus_{i\in I}(X_i,\Sigma_{X_i},\mu_{X_i})$ dari
ukuran-ukuran subruang, jika kita memasangkan $(x,i)\in X'$ dengan
$x\in X$ untuk setiap $i\in I$ dan $x\in X_i$.

\cmmnt{Untuk beberapa sifat elementer (terus terang, saya tidak
mengetahui sifat apa pun yang tidak elementer) dari jumlah langsung,
lihat 214M dan 214Xh-214Xk.}
% 214Xh, 214Xi, 214Xj, 214Xk

\leader{214M}{Proposisi}
Misalkan $\familyiI{(X_i,\Sigma_i,\mu_i)}$ suatu keluarga ruang ukur,
dengan jumlah langsung $(X,\Sigma,\mu)$.   Misalkan $f$ suatu fungsi
bernilai real yang didefinisikan pada suatu himpunan bagian dari $X$.
Untuk setiap $i\in I$, tetapkan $f_i(x)=f(x,i)$ kapan pun
$(x,i)\in\dom f$.

(a) $f$ terukur jika dan hanya jika $f_i$ terukur untuk setiap $i\in I$.

(b) Jika $f$ nonnegatif, maka
$\int fd\mu=\sum_{i\in I}\int f_id\mu_i$ apabila salah satu ruas
terdefinisi dalam $[0,\infty]$.

\proof{{\bf (a)} Untuk $a\in\Bbb R$, tetapkan
$F_a=\{(x,i):(x,i)\in\dom f,\,f(x,i)\ge a\}$.   (i) Jika $f$ terukur,
$i\in I$ dan $a\in\Bbb R$, maka terdapat $E\in\Sigma$ sedemikian sehingga
$F_a=E\cap\dom f$; kini

\Centerline{$\{x:f_i(x)\ge a\}
=\dom f_i\cap\{x:(x,i)\in E\}$}

\noindent termasuk dalam aljabar-sigma subruang pada $\dom f_i$ yang
diinduksi oleh $\Sigma_i$.   Karena $a$ sembarang, $f_i$ terukur.
(ii) Jika setiap $f_i$ terukur dan $a\in\Bbb R$, maka untuk setiap
$i\in I$ terdapat $E_i\in\Sigma_i$ sedemikian sehingga
$\{x:(x,i)\in F_a\}=E_i\cap\dom f_i$; dengan menetapkan
$E=\{(x,i):i\in I,\,x\in E_i\}$, himpunan $F_a=\dom f\cap E$ termasuk
dalam aljabar-sigma subruang pada $\dom f$.   Karena $a$ sembarang,
$f$ terukur.

\medskip

{\bf (b)(i)} Mula-mula andaikan $f$ terukur dan terdefinisi di seluruh
ruang.   Tetapkan
$F_{nk}=\{(x,i):(x,i)\in X,\,f(x,i)\ge 2^{-n}k\}$ untuk $k$,
$n\in\Bbb N$, $g_n=\sum_{k=1}^{4^n}2^{-n}\chi F_{nk}$ untuk
$n\in\Bbb N$, $F_{nki}=\{x:(x,i)\in F_{nk}\}$ untuk $k$, $n\in\Bbb N$
dan $i\in I$, serta $g_{ni}(x)=g_n(x,i)$ untuk $i\in I$, $x\in X_i$.
Maka

$$\eqalign{\int fd\mu
&=\lim_{n\to\infty}\int g_nd\mu
=\sup_{n\in\Bbb N}\sum_{k=1}^{4^n}2^{-n}\mu F_{nk}\cr
&=\sup_{n\in\Bbb N}\sum_{k=1}^{4^n}\sum_{i\in I}2^{-n}\mu_i F_{nki}
=\sum_{i\in I}\sup_{n\in\Bbb N}\sum_{k=1}^{4^n}2^{-n}\mu_i F_{nki}\cr
&=\sum_{i\in I}\sup_{n\in\Bbb N}\int g_{ni}d\mu_i
=\sum_{i\in I}\int f_id\mu_i.\cr}$$

\medskip

\quad{\bf (ii)} Secara umum, jika $\int fd\mu$ terdefinisi, terdapat
suatu fungsi terukur $g:X\to\coint{0,\infty}$ dan suatu himpunan terukur
koterabaikan $E\subseteq\dom f$ sedemikian sehingga $g=f$ pada $E$.
Kini $E_i=\{x:(x,i)\in E\}$ termasuk dalam $\Sigma_i$ untuk setiap
$i$, dan $\sum_{i\in I}\mu_i(X_i\setminus E_i)=\mu(X\setminus E)=0$,
sehingga $E_i$ koterabaikan terhadap $\mu_i$ untuk setiap $i$.   Dengan
menetapkan $g_i(x)=g(x,i)$ untuk $x\in X_i$, (i) memberi tahu kita bahwa

\Centerline{$\sum_{i\in I}\int f_id\mu_i
=\sum_{i\in I}\int g_id\mu_i=\int g\,d\mu=\int fd\mu$.}

\medskip

\quad{\bf (iii)} Sebaliknya, jika $\int f_id\mu_i$ terdefinisi untuk
setiap $i\in I$, maka untuk setiap $i\in I$ kita dapat menemukan fungsi
terukur $g_i:X_i\to\coint{0,\infty}$ dan himpunan terukur koterabaikan
terhadap $\mu_i$, $E_i\subseteq\dom f_i$, sedemikian sehingga $g_i=f_i$
pada $E_i$.   Dengan menetapkan $g(x,i)=g_i(x)$ untuk $i\in I$,
$x\in X_i$, (a) memberi tahu kita bahwa $g$ terukur, sedangkan $g=f$
pada $\{(x,i):i\in I,\,x\in E_i\}$, yang koterabaikan (berdasarkan
perhitungan dalam (ii) tepat di atas); jadi

\Centerline{$\int fd\mu=\int g\,d\mu
=\sum_{i\in I}\int g_id\mu_i=\sum_{i\in I}\int f_id\mu_i$,}

\noindent dengan kembali menggunakan (i) untuk langkah tengah.
}%end of proof of 214M

\leader{214N}{Korolari}
Misalkan $(X,\Sigma,\mu)$ suatu ruang ukur dengan dekomposisi
$\familyiI{X_i}$.   Jika $f$ suatu fungsi bernilai real yang
didefinisikan pada suatu himpunan bagian dari $X$, maka

(a) $f$ terukur jika dan hanya jika $f\restr X_i$ terukur untuk setiap
$i\in I$,

(b) jika $f\ge 0$, maka $\int f=\sum_{i\in I}\int_{X_i}f$ apabila salah
satu ruas terdefinisi dalam $[0,\infty]$.

\proof{ Terapkan 214M pada jumlah langsung dari
$\familyiI{(X_i,\Sigma_{X_i},\mu_{X_i})}$, yang diidentifikasi dengan
$(X,\Sigma,\mu)$ sebagaimana dalam 214L.
}%end of proof of 214N

\leader{*214O}{}\cmmnt{ Saya menyediakan ruang di sini untuk suatu
teorema umum yang menuntut cukup banyak dari pembaca.   Karena itu, saya
perlu mengatakan bahwa saya menyarankan agar teorema ini dilewati pada
pembacaan pertama.   Teorema ini tidak akan dikutip dalam jilid ini;
dalam bentuk lengkapnya di sini, saya tidak memperkirakan akan
menggunakannya di bagian mana pun dalam risalah ini; hanya kasus khusus
214Xm yang cukup sering diterapkan; dan buktinya bergantung pada suatu
konsep (`ideal himpunan') serta suatu teknik (`induksi transfinit', bagian
(d) dari bukti 214P) yang tidak digunakan di tempat lain dalam jilid ini.
Meskipun demikian, `perluasan ukuran' merupakan salah satu tema utama
Jilid 4, dan hasil ini dapat membantu memperjelas beberapa pola yang akan
muncul di sana.

\medskip

\noindent}{\bf Lema} Misalkan $(X,\Sigma,\mu)$ suatu ruang ukur, dan
$\Cal I$ suatu ideal dari himpunan-himpunan bagian $X$, yakni suatu
keluarga himpunan bagian $X$ sedemikian sehingga $\emptyset\in\Cal I$,
$I\cup J\in\Cal I$ untuk semua $I$, $J\in\Cal I$, dan $I\in\Cal I$
setiap kali $I\subseteq J\in\Cal I$.   Maka terdapat suatu ukuran
$\lambda$ pada $X$ sedemikian sehingga
$\Sigma\cup\Cal I\subseteq\dom\lambda$,
$\mu E=\lambda E+\sup_{I\in\Cal I}\mu^*(E\cap I)$ untuk setiap
$E\in\Sigma$, dan $\lambda I=0$ untuk setiap $I\in\Cal I$.

\proof{{\bf (a)} Misalkan $\Lambda$ himpunan semua $F\subseteq X$
sedemikian sehingga terdapat $E\in\Sigma$ dan suatu
$\Cal J\subseteq\Cal I$ terhitung dengan
$E\symmdiff F\subseteq\bigcup\Cal J$.   Maka $\Lambda$ merupakan suatu
aljabar-sigma dari himpunan-himpunan bagian $X$ yang memuat
$\Sigma\cup\Cal I$.   \Prf\ $\Sigma\subseteq\Lambda$ karena
$E\symmdiff E\subseteq\bigcup\emptyset$ untuk setiap $E\in\Sigma$.
$\Cal I\subseteq\Lambda$ karena
$\emptyset\symmdiff I\subseteq\bigcup\{I\}$ untuk setiap $I\in\Cal I$.
Khususnya, $\emptyset\in\Lambda$.   Jika $F\in\Lambda$, misalkan
$E\in\Sigma$ dan $\Cal J\subseteq\Cal I$ sedemikian sehingga $\Cal J$
terhitung dan $F\symmdiff E\subseteq\bigcup\Cal J$; maka
$(X\setminus F)\symmdiff(X\setminus E)\subseteq\bigcup\Cal J$, sehingga
$X\setminus F\in\Lambda$.   Jika $\sequencen{F_n}$ suatu barisan dalam
$\Lambda$ dengan gabungan $F$, maka untuk setiap $n\in\Bbb N$ pilih
$E_n\in\Sigma$, $\Cal J_n\subseteq\Cal I$ sedemikian sehingga
$\Cal J_n$ terhitung dan
$E_n\symmdiff F_n\subseteq\bigcup\Cal J_n$; maka
$E=\bigcup_{n\in\Bbb N}E_n$ termasuk dalam $\Sigma$,
$\Cal J=\bigcup_{n\in\Bbb N}\Cal J_n$ merupakan himpunan bagian
terhitung dari $\Cal I$, dan $E\symmdiff F\subseteq\bigcup\Cal J$,
sehingga $F\in\Lambda$.   Jadi $\Lambda$ suatu aljabar-sigma.\ \Qed

\medskip

{\bf (b)} Untuk $F\in\Lambda$ tetapkan

\Centerline{$\lambda F=\sup\{\mu E:E\in\Sigma$,
$E\subseteq F$, $\mu^*(E\cap I)=0$ untuk setiap $I\in\Cal I\}$.}

\noindent Maka $\lambda$ suatu ukuran.   \Prf\ Satu-satunya himpunan
bagian dari $\emptyset$ adalah $\emptyset$, sehingga
$\lambda\emptyset=0$.   Misalkan $\sequencen{F_n}$ suatu barisan saling
lepas dalam $\Lambda$ dengan gabungan $F$; tetapkan
$u=\sum_{n=0}^{\infty}\lambda F_n$.
(i) Jika $E\in\Sigma$, $E\subseteq F$, dan $\mu^*(E\cap I)=0$ untuk
setiap $I\in\Cal I$, maka untuk setiap $n$ tetapkan $E_n=E\cap F_n$.
Karena $\mu^*(E_n\cap I)=0$ untuk setiap $I\in\Cal I$,
$\mu E_n\le\lambda F_n$ untuk setiap $n$.   Kini
$\sequencen{E_n}$ saling lepas dan mempunyai gabungan $E$, sehingga

\Centerline{$\mu E
=\sum_{n=0}^{\infty}\mu E_n\le\sum_{n=0}^{\infty}\lambda F_n=u$.}

\noindent Karena $E$ sembarang, $\lambda F\le u$.   (ii) Ambil sembarang
$\gamma<u$.   Untuk $n\in\Bbb N$, tetapkan
$\gamma_n=\lambda F_n-2^{-n-1}\min(1,u-\gamma)$ jika $\lambda F_n$
berhingga, dan $\gamma$ jika tidak.   Untuk setiap $n$, kita dapat
menemukan $E_n\in\Sigma$ sedemikian sehingga $E_n\subseteq F_n$,
$\mu^*(E_n\cap I)=0$ untuk setiap $I\in\Cal I$, dan
$\mu E_n\ge\gamma_n$.   Tetapkan $E=\bigcup_{n\in\Bbb N}E_n$; maka
$E\subseteq F$ dan $E\cap I=\bigcup_{n\in\Bbb N}E_n\cap I$ terabaikan
terhadap $\mu$ untuk setiap $I\in\Cal I$, sehingga
$\lambda F\ge\mu E=\sum_{n=0}^{\infty}\mu E_n\ge\gamma$. Karena
$\gamma$ sembarang, $\lambda F\ge u$.   (iii) Karena
$\sequencen{F_n}$ sembarang, $\lambda$ suatu ukuran.\ \Qed

\medskip

{\bf (c)} Sekarang ambil sembarang $E\in\Sigma$ dan tetapkan
$u=\sup_{I\in\Cal I}\mu^*(E\cap I)$.   Jika $u=\infty$, tentu saja kita
mempunyai $\mu E=\infty=\lambda E+u$. Jika tidak, misalkan
$\sequencen{I_n}$ suatu barisan dalam $\Cal I$ sedemikian sehingga
$\lim_{n\to\infty}\mu^*(E\cap I_n)=u$; dengan mengganti $I_n$ oleh
$\bigcup_{m\le n}I_m$ untuk setiap $n$ jika perlu, kita dapat mengandaikan
$\sequencen{I_n}$ tak-menurun.   Tetapkan
$A=E\cap\bigcup_{n\in\Bbb N}I_n$; karena $E\cap I_n$ mempunyai ukuran
luar berhingga untuk setiap $n$, $A$ dapat ditutupi oleh suatu barisan
himpunan berukuran hingga dan mempunyai selubung terukur $H$ terhadap
$\mu$ yang termuat dalam $E$ (132Ee).   Perhatikan bahwa

\Centerline{$\mu H=\mu^*A=\sup_{n\in\Bbb N}\mu^*(E\cap I_n)=u$}

\noindent berdasarkan 132Ae.

Tetapkan $G=E\setminus H$.   Maka $\mu^*(G\cap I)=0$ untuk setiap
$I\in\Cal I$. \Prf\ Untuk sembarang $n\in\Bbb N$ terdapat
$F\in\Sigma$ sedemikian sehingga $F\supseteq E\cap(I_n\cup I)$ dan
$\mu F\le u$; dalam hal ini

\Centerline{$\mu^*(G\cap I)+\mu^*(E\cap I_n)
\le\mu(F\setminus H)+\mu(F\cap H)\le u$.}

\noindent Karena $n$ sembarang, $\mu^*(G\cap I)=0$.\ \QeD\   Oleh karena itu

\Centerline{$u+\lambda E\ge\mu H+\mu G=\mu E$.}

\noindent Di pihak lain, jika $F\in\Sigma$ sedemikian sehingga
$F\subseteq E$ dan $\mu^*(F\cap I)=0$ untuk setiap $I\in\Cal I$, maka

\Centerline{$\mu^*(E\cap I_n)
\le\mu(E\setminus F)+\mu^*(F\cap I_n)=\mu(E\setminus F)$}

\noindent untuk setiap $n$, sehingga

\Centerline{$u+\mu F\le\mu(E\setminus F)+\mu F=\mu E$;}

\noindent karena $F$ sembarang, $u+\lambda E\le\mu E$.

\medskip

{\bf (d)} Jika $J\in\Cal I$, $F\in\Sigma$, $F\subseteq J$, dan
$\mu^*(F\cap I)=0$ untuk setiap $I\in\Cal I$, maka
$F\cap J=F$ terabaikan terhadap $\mu$; karena $F$ sembarang,
$\lambda J=0$. Jadi $\lambda$ mempunyai semua sifat yang diperlukan.
}%end of proof of 214O

\leader{*214P}{Teorema}
Misalkan $(X,\Sigma,\mu)$ suatu ruang ukur, dan $\Cal A$ suatu keluarga
himpunan bagian dari $X$ yang terurut baik oleh relasi $\subseteq$.
Maka terdapat suatu perluasan $\mu$ menjadi suatu ukuran $\lambda$ pada
$X$ sedemikian sehingga $\lambda(E\cap A)$ terdefinisi dan sama dengan
$\mu^*(E\cap A)$ setiap kali $E\in\Sigma$ dan $A\in\Cal A$.

\proof{{\bf (a)} Dengan menambahkan $\emptyset$ dan $X$ ke $\Cal A$ jika
perlu, kita dapat mengandaikan bahwa $\Cal A$ mempunyai $\emptyset$
sebagai anggota terkecil dan $X$ sebagai anggota terbesar.   Berdasarkan
2A1Dg, $\Cal A$ isomorfik, sebagai himpunan terurut, dengan suatu
ordinal; karena $\Cal A$ mempunyai anggota terbesar, ordinal ini
merupakan ordinal penerus, yang dapat dinyatakan sebagai $\zeta+1$;
misalkan $\xi\mapsto A_{\xi}:\zeta+1\to\Cal A$ isomorfisme urutannya,
sehingga $\langle A_{\xi}\rangle_{\xi\le\zeta}$ merupakan keluarga
himpunan bagian $X$ yang tak-menurun, $A_0=\emptyset$ dan $A_{\zeta}=X$.

\medskip

{\bf (b)} Untuk setiap ordinal $\xi\le\zeta$, tuliskan $\mu_{\xi}$ untuk
ukuran subruang pada $A_{\xi}$, $\Sigma_{\xi}$ untuk domainnya, dan
$\Cal I_{\xi}$ untuk $\bigcup_{\eta<\xi}\Cal PA_{\eta}$. Karena
$A_{\eta}\cup A_{\eta'}=A_{\max(\eta,\eta')}$ untuk $\eta$,
$\eta'<\xi$, $\Cal I_{\xi}$ merupakan ideal dari himpunan-himpunan bagian
$A_{\xi}$.   Berdasarkan 214O, kita mempunyai suatu ukuran
$\lambda_{\xi}$ pada $A_{\xi}$, dengan domain
$\Lambda_{\xi}\supseteq\Sigma_{\xi}\cup\Cal I_{\xi}$, sedemikian sehingga
$\mu_{\xi}E=\lambda_{\xi}E+\sup_{I\in\Cal I_{\xi}}\mu_{\xi}^*(E\cap I)$
untuk setiap $E\in\Sigma_{\xi}$ dan $\lambda_{\xi}I=0$ untuk setiap
$I\in\Cal I_{\xi}$.   Karena setiap anggota $\Cal I_{\xi}$ termuat dalam
$A_{\eta}$ untuk suatu $\eta<\xi$, kita mempunyai

\Centerline{$\mu^*(E\cap A_{\xi})
=\lambda_{\xi}(E\cap A_{\xi})+\sup_{\eta<\xi}\mu_{\xi}^*(E\cap A_{\eta})
=\lambda_{\xi}(E\cap A_{\xi})+\sup_{\eta<\xi}\mu^*(E\cap A_{\eta})$}

\noindent (214Cd) untuk setiap $E\in\Sigma$.   Selain itu, tentu saja,
$\lambda_{\xi}A_{\eta}=0$ untuk setiap $\eta<\xi$.

\medskip

{\bf (c)} Sekarang tetapkan

\Centerline{$\Lambda=\{F:F\subseteq X$, $F\cap A_{\xi}\in\Lambda_{\xi}$ untuk
setiap $\xi\le\zeta\}$,}

\Centerline{$\lambda F=\sum_{\xi\le\zeta}\lambda_{\xi}(F\cap A_{\xi})$}

\noindent untuk setiap $F\in\Lambda$.   Karena $\Lambda_{\xi}$ suatu
aljabar-sigma dari himpunan-himpunan bagian $A_{\xi}$ untuk setiap
$\xi$, $\Lambda$ suatu aljabar-sigma dari himpunan-himpunan bagian
$X$; karena setiap $\lambda_{\xi}$ suatu ukuran, $\lambda$ juga demikian.
Jika $E\in\Sigma$, maka

\Centerline{$E\cap A_{\xi}\in\Sigma_{\xi}\subseteq\Lambda_{\xi}$}

\noindent untuk setiap $\xi$, sehingga $E\in\Lambda$.   Jika
$\eta\le\zeta$, maka untuk setiap $\xi\le\zeta$, berlaku salah satu dari
dua hal: $\eta<\xi$ dan

\Centerline{$A_{\eta}\cap A_{\xi}=A_{\eta}
\in\Cal I_{\xi}\subseteq\Lambda_{\xi}$}

\noindent atau $\eta\ge\xi$ dan $A_{\eta}\cap A_{\xi}=A_{\xi}$ termasuk
dalam $\Lambda_{\xi}$.   Jadi $A_{\eta}\in\Lambda$ untuk setiap
$\eta\le\zeta$.

\medskip

{\bf (d)} Akhirnya,
$\lambda(E\cap A_{\xi})=\mu^*(E\cap A_{\xi})$ setiap kali $E\in\Sigma$
dan $\xi\le\zeta$.   \Prf\Quer\ Jika tidak, karena ordinal $\zeta+1$
terurut baik, terdapat $\xi$ terkecil sedemikian sehingga
$\lambda(E\cap A_{\xi})\ne\mu^*(E\cap A_{\xi})$. Karena
$A_0=\emptyset$, tentu kita mempunyai
$\lambda(E\cap A_0)=\mu^*(E\cap A_0)$ dan $\xi>0$.
Perhatikan bahwa jika $\eta>\xi$, maka
$\lambda_{\eta}(E\cap A_{\xi})=0$; sehingga

\Centerline{$\lambda(E\cap A_{\xi})
=\sum_{\eta\le\xi}\lambda_{\eta}(E\cap A_{\xi}\cap A_{\eta})
=\sum_{\eta\le\xi}\lambda_{\eta}(E\cap A_{\eta})$.}

Kini

$$\eqalignno{\mu^*(E\cap A_{\xi})
&=\lambda_{\xi}(E\cap A_{\xi})+\sup_{\xi'<\xi}\mu^*(E\cap A_{\xi'})\cr
\displaycause{(b) di atas}
&=\lambda_{\xi}(E\cap A_{\xi})
  +\sup_{\xi'<\xi}\sum_{\eta\le\xi'}\lambda_{\eta}(E\cap A_{\eta})\cr
\displaycause{karena $\xi$ merupakan ordinal bermasalah pertama}
&=\lambda_{\xi}(E\cap A_{\xi})
  +\sup_{\xi'<\xi}\sup_{K\subseteq\xi'+1\text{ berhingga}}
  \sum_{\eta\in K}\lambda_{\eta}(E\cap A_{\eta})\cr
\displaycause{lihat definisi `jumlah' dalam 112Bd, atau 226A di bawah}
&=\lambda_{\xi}(E\cap A_{\xi})
  +\sup_{K\subseteq\xi\text{ berhingga}}
  \sum_{\eta\in K}\lambda_{\eta}(E\cap A_{\eta})\cr
&=\sup_{K\subseteq\xi+1\text{ berhingga}}
  \sum_{\eta\in K}\lambda_{\eta}(E\cap A_{\eta})
=\sum_{\eta\le\xi}\lambda_{\eta}(E\cap A_{\eta})
\ne\mu^*(E\cap A_{\xi})\cr}$$

\noindent menurut pemilihan $\xi$; tetapi ini mustahil.\ \Bang\Qed

Khususnya,

\Centerline{$\lambda E=\lambda(E\cap A_{\zeta})=\mu^*(E\cap A_{\zeta})
=\mu E$}

\noindent untuk setiap $E\in\Sigma$.
Ini menyelesaikan bukti teorema.
}%end of proof of 214P

\leader{*214Q}{Proposisi}\dvAformerly{2{}14Xd}
Andaikan $(X,\Sigma,\mu)$ suatu ruang ukur tanpa atom dan $Y$ suatu
himpunan bagian dari $X$ sedemikian sehingga ukuran subruang $\mu_Y$
semihingga.   Maka $\mu_Y$ tanpa atom.
%for 548B 548C
%\query what about general $Y$?

\proof{ Misalkan $F\subseteq Y$ sedemikian sehingga $\mu_YF$ terdefinisi
dan tidak sama dengan $0$. Karena $\mu_Y$ semihingga, terdapat
$F'\subseteq F$ sedemikian sehingga $\mu_YF'$ terdefinisi, berhingga, dan
tidak nol.   Dalam hal ini, $\mu^*F'=\mu_YF'$ berhingga, sehingga $F'$
mempunyai suatu selubung terukur $E$, katakanlah, terhadap $\mu$.
Karena $\mu$ tanpa atom, terdapat $E_1\in\Sigma$ sedemikian sehingga
$E_1\subseteq E$ dan baik $E_1$ maupun $E\setminus E_1$ tidak
terabaikan terhadap $\mu$.   Kini $F\cap E_1$ diukur oleh $\mu_Y$ dan

\Centerline{$\mu_Y(F\cap E_1)\ge\mu^*(F'\cap E_1)=\mu(E\cap E_1)>0$,}

\Centerline{$\mu_Y(F\setminus E_1)\ge\mu^*(F'\setminus E_1)
=\mu(E\setminus E_1)>0$.}

\noindent Karena $F$ sembarang, $\mu_Y$ tanpa atom.
}%end of proof of 214Q

\exercises{
\leader{214X}{Latihan dasar (a)}
%\spheader 214Xa
Misalkan $(X,\Sigma,\mu)$ suatu ruang ukur yang dapat dilokalkan.
Tunjukkan bahwa terdapat $E\in\Sigma$ sedemikian sehingga ukuran subruang
$\mu_E$ atomik murni dan $\mu_{X\setminus E}$ tanpa atom.
%214/

\spheader 214Xb Misalkan $X$ suatu himpunan, $\theta$ suatu ukuran luar
reguler pada $X$, dan $Y$ suatu himpunan bagian dari $X$.   Misalkan
$\mu$ ukuran pada $X$ yang didefinisikan dengan metode \Caratheodory\
dari $\theta$, $\mu_Y$ ukuran subruang pada $Y$, dan $\nu$ ukuran pada
$Y$ yang didefinisikan dengan metode \Caratheodory\ dari
$\theta\restrp\Cal PY$.   Tunjukkan bahwa jika $\mu_Y$ ditentukan secara
lokal (khususnya, jika $\mu$ ditentukan secara lokal dan dapat
dapat dilokalkan), maka $\nu=\mu_Y$.
%214I

\spheader 214Xc Misalkan $(X,\Sigma,\mu)$ suatu ruang ukur yang dapat
dapat dilokalkan, dan $Y$ suatu himpunan bagian dari $X$ sedemikian sehingga
ukuran subruang $\mu_Y$ semihingga.   Tunjukkan bahwa $\mu_Y$ dapat
dapat dilokalkan.
%214I

\sqheader 214Xd Misalkan $(X,\Sigma,\mu)$ suatu ruang ukur, dan $Y$
suatu himpunan bagian dari $X$ sedemikian sehingga ukuran subruang
$\mu_Y$ semihingga. (i) Tunjukkan bahwa suatu himpunan $F\subseteq Y$
merupakan atom bagi $\mu_Y$ jika dan hanya jika berbentuk $E\cap Y$
dengan $E$ suatu atom bagi $\mu$. (ii) Tunjukkan bahwa jika $\mu$ murni
atomik, maka $\mu_Y$ juga demikian.
%214I  used in 5x47A

\spheader 214Xe Misalkan $(X,\Sigma,\mu)$ suatu ruang ukur yang dapat
dapat dilokalkan, dan $Y$ sembarang himpunan bagian dari $X$.   Tunjukkan
bahwa versi c.l.d.\ dari ukuran subruang pada $Y$ dapat dilokalkan.
%214I

\spheader 214Xf Misalkan $(X,\Sigma,\mu)$ suatu ruang ukur dengan
himpunan-himpunan terabaikan yang ditentukan secara lokal, dan $Y$ suatu
himpunan bagian dari $X$, dengan ukuran subruangnya $\mu_Y$. Tunjukkan
bahwa $\mu_Y$ mempunyai himpunan-himpunan terabaikan yang ditentukan
secara lokal.
%214I

\sqheader 214Xg Misalkan $(X,\Sigma,\mu)$ suatu ruang ukur.   Tunjukkan
bahwa $(X,\Sigma,\mu)$ mempunyai himpunan-himpunan terabaikan yang
ditentukan secara lokal jika dan hanya jika ukuran subruang $\mu_Y$
semihingga untuk setiap $Y\subseteq X$.
%214I, 214Xf

\sqheader 214Xh Misalkan $\langle (X_i,\Sigma_i,\mu_i)\rangle_{i\in I}$
suatu keluarga ruang ukur, dengan jumlah langsung $(X,\Sigma,\mu)$
(214L).   Tetapkan $X'_i=X_i\times\{i\}\subseteq X$ untuk setiap
$i\in I$.   Tunjukkan bahwa $X'_i$, dengan ukuran subruangnya, isomorfik
dengan $(X_i,\Sigma_i,\mu_i)$.   Dalam keadaan apa
$\langle X'_i\rangle_{i\in I}$ merupakan dekomposisi $X$?   Tunjukkan
bahwa $\mu$ lengkap, atau dapat dilokalkan secara ketat, atau dapat
dilokalkan, atau ditentukan secara lokal, atau semihingga, atau tanpa
atom, %used in 548C
atau atomik murni, jika dan hanya jika setiap $\mu_i$ demikian.
Tunjukkan bahwa suatu ruang ukur dapat dilokalkan secara ketat jika dan
hanya jika ruang tersebut isomorfik dengan suatu jumlah langsung dari
ruang-ruang berhingga total.
%214L

\sqheader 214Xi Misalkan
$\langle(X_i,\Sigma_i,\mu_i)\rangle_{i\in I}$ suatu keluarga ruang ukur,
dan $(X,\Sigma,\mu)$ jumlah langsungnya.   Tunjukkan bahwa pelengkapan
$(X,\Sigma,\mu)$ dapat diidentifikasi dengan jumlah langsung dari
pelengkapan ruang-ruang $(X_i,\Sigma_i,\mu_i)$, dan bahwa versi c.l.d.\
dari $(X,\Sigma,\mu)$ dapat diidentifikasi dengan jumlah langsung dari
versi-versi c.l.d.\ ruang $(X_i,\Sigma_i,\mu_i)$.
%214L

\spheader 214Xj Misalkan
$\langle(X_i,\Sigma_i,\mu_i)\rangle_{i\in I}$ suatu keluarga ruang ukur.
Tunjukkan bahwa jumlah langsungnya mempunyai himpunan-himpunan terabaikan
yang ditentukan secara lokal jika dan hanya jika setiap $\mu_i$
mempunyainya.
%214L, 214Xf

\spheader 214Xk
Misalkan $\langle(X_i,\Sigma_i,\mu_i)\rangle_{i\in I}$ suatu keluarga
ruang ukur, dan $(X,\Sigma,\mu)$ jumlah langsungnya.   Tunjukkan bahwa
$(X,\Sigma,\mu)$ mempunyai sifat selubung terukur (213Xl) jika dan hanya
jika setiap $(X_i,\Sigma_i,\mu_i)$ mempunyainya.
%214L

\spheader 214Xl
Misalkan $(X,\Sigma,\mu)$ suatu ruang ukur, $Y$ suatu himpunan bagian dari
$X$, dan $f:X\to[0,\infty]$ suatu fungsi sedemikian sehingga $\int_Yf$
terdefinisi dalam $[0,\infty]$. Tunjukkan bahwa
$\int_Yf=\overline{\int}f\times\chi Y\,d\mu$.
%214E

\sqheader 214Xm\dvAnew{2009}
Tuliskan bukti langsung untuk 214P dalam kasus khusus ketika
$\Cal A=\{A\}$.   ({\it Petunjuk\/}: untuk $E$, $F\in\Sigma$,

\Centerline{$\lambda((E\cap A)\cup(F\setminus A))
=\mu^*(E\cap A)+\sup\{\mu G:G\in\Sigma$, $G\subseteq F\setminus A\}$.)}
%214P

\sqheader 214Xn\dvAnew{2009}
Misalkan $(X,\Sigma,\mu)$ suatu ruang ukur dan $\Cal A$ suatu keluarga
berhingga dari himpunan-himpunan bagian $X$.   Tunjukkan bahwa terdapat
suatu ukuran pada $X$, yang memperluas $\mu$, dan mengukur setiap anggota
$\Cal A$.
%214Xm 214P

\leader{214Y}{Latihan lanjutan (a)}\dvAnew{2009}
%\spheader 214Ya
Misalkan $(X,\Sigma,\mu)$ suatu ruang ukur dan $A$ suatu himpunan bagian
dari $X$ sedemikian sehingga ukuran subruang pada $A$ semihingga.
Tetapkan $\alpha=\sup\{\mu E:E\in\Sigma$, $E\subseteq A\}$.
Tunjukkan bahwa jika $\alpha\le\gamma\le\mu^*A$, maka terdapat suatu
ukuran $\lambda$ pada $X$, yang memperluas $\mu$, sedemikian sehingga
$\lambda A=\gamma$.
%214P 214Xm

\spheader 214Yb\dvAnew{2009}
Misalkan $(X,\Sigma,\mu)$ suatu ruang ukur dan
$\family{n}{\Bbb Z}{A_n}$ suatu barisan dua-arah dari himpunan-himpunan
bagian $X$ sedemikian sehingga $A_m\subseteq A_n$ setiap kali $m\le n$
dalam $\Bbb Z$.   Tunjukkan bahwa terdapat suatu ukuran pada $X$, yang
memperluas $\mu$, dan mengukur setiap $A_n$.
\Hint{gunakan 214P dua kali.}
%214P

\spheader 214Yc\dvAnew{2009}
Misalkan $X$ suatu himpunan dan $\Cal A$ suatu keluarga himpunan bagian
dari $X$. Tunjukkan bahwa pernyataan-pernyataan berikut ekuivalen:
(i) untuk setiap ukuran $\mu$ pada $X$, terdapat suatu ukuran pada $X$
yang memperluas $\mu$ dan mengukur setiap anggota $\Cal A$; (ii) untuk
setiap ukuran berhingga total $\mu$ pada $X$, terdapat suatu ukuran pada
$X$ yang memperluas $\mu$ dan mengukur setiap anggota $\Cal A$.
\Hint{213Xa.}

\spheader 214Yd\dvAnew{2015}
Untuk latihan ini saja, saya akan mengatakan bahwa suatu ukuran $\mu$
pada suatu himpunan $X$ bersifat {\bf tak pernah mengukur semua} jika,
setiap kali $A\subseteq X$ tidak terabaikan terhadap $\mu$, terdapat
suatu himpunan bagian dari $A$ yang tidak diukur oleh ukuran subruang
pada $A$.   Tunjukkan bahwa jika $X$ suatu himpunan dan
$\mu_0,\ldots,\mu_n$ merupakan ukuran lengkap berhingga total yang tak
pernah mengukur semua pada $X$, maka terdapat himpunan-himpunan saling
lepas $A_0,\ldots,A_n\subseteq X$ sedemikian sehingga
$\mu_i^*A_i=\mu_iX$ untuk setiap $i\le n$.   \Hint{mulailah dengan kasus
$n=1$, $\mu_0=\mu_1$.}
%mt21bits
}%end of exercises

\endnotes{
\Notesheader{214} Saya menguraikan bagian pertama dari bagian ini,
hingga 214H, secara perlahan dan cermat, karena meskipun tidak satu pun
argumennya mendalam (214Eb merupakan yang terpanjang), pola yang dibentuk
oleh hasil-hasil tersebut tidak selalu mudah diperkirakan.   Terdapat
contoh tandingan terhadap perluasan 214H/214Xb yang tampak menggiurkan
dalam 216Xb.

Pesan dari bagian kedua dari bagian ini (214I-214L, 214Q) adalah bahwa
subruang mewarisi banyak, tetapi tidak semua, sifat suatu ruang ukur; dan
khususnya dapat timbul kesulitan dengan sifat semihingga, kecuali jika kita
mempunyai himpunan-himpunan terabaikan yang ditentukan secara lokal
(214Xg).   (Saya memberikan contoh dalam 216Xa.)   Tentu saja 213Hb
menunjukkan bahwa jika kita mulai dengan suatu ruang yang dapat
dapat dilokalkan, kita dapat mengubahnya menjadi ruang yang lengkap,
ditentukan secara lokal, dan dapat dilokalkan tanpa terlalu merusak
struktur ruang tersebut, sehingga kesulitannya biasanya dapat diatasi.

Kasus 214P yang paling penting adalah ketika $\Cal A=\{A\}$ suatu
himpunan tunggal, sehingga argumennya menjadi jauh lebih sederhana (214Xm).
Dalam \S439 dari Jilid 4 saya akan kembali pada masalah memperluas suatu
ukuran ke aljabar-sigma lebih besar yang diberikan, tanpa adanya
struktur bantu yang berguna.   Bagian tersebut terutama akan menawarkan
contoh-contoh tandingan, khususnya menunjukkan bahwa tidak ada teorema
umum yang memperluas 214Xn dari keluarga berhingga ke keluarga terhitung,
dan bahwa syarat-syarat khusus dalam 214P dan 214Yb ada dengan alasan
yang baik.   Namun, dalam Bab 54 dan \S552 dari Jilid 5 saya akan membahas
sistem-sistem matematika yang memungkinkan teorema perluasan jauh lebih
kuat, setidaknya jika kita mulai dari ukuran Lebesgue.
}%end of notes

\discrpage


